% Compile: xelatex CPE342_Assignment 3_main.tex (รัน 2 รอบ)
\documentclass[12pt,a4paper]{article}

\usepackage{fontspec}
\usepackage{polyglossia}
\setdefaultlanguage{thai}
\setotherlanguage{english}

\setmainfont{Sarabun}[
    BoldFont       = Sarabun-Bold,
    ItalicFont     = Sarabun-Italic,
    BoldItalicFont = Sarabun-BoldItalic,
]
\setmonofont{Consolas}[Scale=0.92]
\newfontfamily\thaifonttt{Consolas}[Scale=0.92]

\usepackage{amsmath, amssymb, mathtools}
\usepackage{graphicx}
\usepackage{booktabs}
\usepackage{geometry}
\usepackage{xcolor}
\usepackage{float}
\usepackage{enumitem}
\usepackage{tcolorbox}
\usepackage{needspace}
\usepackage{caption}
\usepackage{listings}
\usepackage{upquote}
\tcbuselibrary{skins, breakable}

\XeTeXlinebreaklocale "th"
\XeTeXlinebreakskip = 0pt plus 1pt
\sloppy

\definecolor{captioncol}{RGB}{80, 80, 80}
\definecolor{accentblue}{RGB}{37, 99, 235}
\definecolor{codebg}{RGB}{248, 250, 252}
\definecolor{codeframe}{RGB}{203, 213, 225}
\definecolor{codeaccent}{HTML}{0969DA}
\newcommand{\codevar}[1]{\textcolor{codeaccent}{\textbf{#1}}}
\definecolor{keyword}{RGB}{124, 58, 237}
\definecolor{string}{RGB}{22, 163, 74}
\definecolor{comment}{RGB}{100, 116, 139}
\definecolor{output}{RGB}{30, 58, 138}
\captionsetup{font={small, color=captioncol}, justification=centering}

\lstdefinestyle{pycode}{
    language=Python,
    backgroundcolor=\color{codebg},
    basicstyle=\ttfamily\small,
    keywordstyle=\color{keyword}\bfseries,
    stringstyle=\color{string},
    commentstyle=\color{comment}\itshape,
    numberstyle=\tiny\color{comment},
    numbers=left,
    numbersep=8pt,
    frame=single,
    rulecolor=\color{codeframe},
    breaklines=true,
    breakatwhitespace=false,
    showstringspaces=false,
    tabsize=4,
    xleftmargin=16pt,
    xrightmargin=4pt,
    columns=flexible,
    keepspaces=true,
    literate={≈}{{$\approx$}}1 {α}{{$\alpha$}}1 {β}{{$\beta$}}1 {λ}{{$\lambda$}}1 {Λ}{{$\Lambda$}}1 {χ}{{$\chi$}}1
}

\lstdefinestyle{output}{
    basicstyle=\ttfamily\small\color{output},
    backgroundcolor=\color{codebg},
    frame=single,
    rulecolor=\color{codeframe},
    numbers=none,
    xleftmargin=8pt,
    xrightmargin=4pt,
    breaklines=true,
    breakatwhitespace=false,
    columns=flexible,
    keepspaces=true,
    literate={≈}{{$\approx$}}1 {α}{{$\alpha$}}1 {β}{{$\beta$}}1 {λ}{{$\lambda$}}1 {Λ}{{$\Lambda$}}1 {χ}{{$\chi$}}1
}

\widowpenalty=10000
\clubpenalty=10000

\geometry{top=2.5cm, bottom=2.5cm, left=2.5cm, right=2.5cm, headheight=14pt}

\newtcolorbox{ansbox}[1][]{
    colback=green!5!white, colframe=green!50!black,
    boxrule=0.5pt, arc=4pt, left=8pt, right=8pt, top=5pt, bottom=5pt, #1
}

% =================================================================
\begin{document}

\begin{center}
    {\Large\textbf{CPE 342 Machine Learning}}\\[2pt]
    {\Large\textbf{Assignment 3: Regression — Survival Analysis}}\\[4pt]
    {\large Estimating Customer Lifetime on Telco Churn Dataset}\\[6pt]
    \textcolor{accentblue}{67070501042 วิศิษฐ์ สุวรรณเนาว์}
\end{center}

\noindent\rule{\textwidth}{0.4pt}\vspace{4pt}

\section*{1. ภาพรวมของปัญหาและวัตถุประสงค์\\ \large(Problem Overview \& Objectives)}

ในงานมอบหมายส่วนของ Activity นี้ เราศึกษาและวิเคราะห์พฤติกรรมการคงอยู่และการยกเลิกบริการของลูกค้า (Customer Retention \& Churn Dynamics) โดยประยุกต์ใช้ระเบียบวิธี \textbf{การวิเคราะห์การอยู่รอด (Survival Analysis)} ด้วยตัวประมาณค่าแบบนอนพารามิเตอร์ \textbf{Kaplan-Meier Estimator} บนชุดข้อมูล \codevar{Telco-Churn.csv} ซึ่งประกอบด้วยข้อมูลลูกค้ารวม $N = 7,043$ ราย

\subsection*{วัตถุประสงค์หลักของการวิเคราะห์}
\begin{enumerate}[leftmargin=18pt, itemsep=2pt]
    \item เพื่อสร้างตัวแปรระยะเวลาการใช้บริการ $T = \text{tenure}$ (หน่วย: เดือน) และตัวบ่งชี้การเกิดเหตุการณ์ยกเลิกบริการ $E = \mathbb{I}(\text{Churn} == \text{'Yes'})$
    \item เพื่อจัดการปัญหา \textbf{ข้อมูลถูกจำกัดช่วงทางขวา (Right-Censoring)} สำหรับลูกค้าที่ยังคงใช้บริการอยู่ ณ ปัจจุบัน ($E = 0$)
    \item เพื่อประมาณค่าฟังก์ชันการอยู่รอดของประชากรโดยรวม $\hat{S}(t)$ พร้อมคำนวณช่วงความเชื่อมั่น 95\% ด้วยสูตรของ Greenwood
    \item เพื่อเปรียบเทียบวิถีการอยู่รอดและความเสี่ยงในการยกเลิกบริการระหว่างช่องทางการชำระเงินทั้ง 4 รูปแบบ (\codevar{PaymentMethod}): \textbf{Electronic check}, \textbf{Mailed check}, \textbf{Bank transfer (automatic)} และ \textbf{Credit card (automatic)}
    \item เพื่อทดสอบนัยสำคัญทางสถิติด้วยการทดสอบ \textbf{Log-Rank Test} (ทั้งแบบ Multivariate และ Pairwise) เพื่อยืนยันความแตกต่างของเส้นโค้งการอยู่รอดอย่างมีนัยสำคัญ ($p < 0.001$)
    \item เพื่อนำผลลัพธ์เชิงปริมาณและค่าความน่าจะเป็น ณ ช่วงเวลาสำคัญ ($t = 12, 24, 36, 48, 60, 72$ เดือน) ไปแปลผลเชิงธุรกิจและเสนอแนะแนวทางลดอัตราการยกเลิกบริการ (Churn Mitigation Strategy)
\end{enumerate}

\section*{2. กรอบทฤษฎีและคณิตศาสตร์\\ \large(Theoretical Framework \& Mathematical Formulations)}

\subsection*{2.1 ข้อจำกัดของ Regression ทั่วไปกับข้อมูล Time-to-Event\\ \normalsize(Limitations of Standard Regression on Censored Data)}
ในการวิเคราะห์ข้อมูลระยะเวลาจนกระทั่งเกิดเหตุการณ์ (Time-to-Event Data) เช่น การยกเลิกบริการของลูกค้า การใช้สมการถดถอยเชิงเส้น (Linear Regression) หรือ Logistic Regression ทั่วไปมีข้อจำกัดร้ายแรงดังนี้ (อิงตาม Lecture 4 หน้า 27--33):
\begin{itemize}[leftmargin=18pt, itemsep=2pt]
    \item \textbf{การติดปัญหา Right-Censoring}: ลูกค้าส่วนใหญ่ ($73.46\%$) ยังไม่ยกเลิกบริการในวันที่เก็บข้อมูล เรารู้เพียงว่าระยะเวลาของลูกค้าเหล่านี้มีค่าอย่างน้อยเท่ากับ $t$ ($T \ge t$) แต่ไม่รู้เวลาสิ้นสุดจริง หากตัดข้อมูลเหล่านี้ทิ้งจะเกิดภาวะความลำเอียงจากการคัดออก (Selection Bias) หรือหากนับว่า $t$ คือเวลาสิ้นสุด โมเดลจะประเมินอายุขัยของลูกค้าต่ำกว่าความเป็นจริง (Underestimation)
    \item \textbf{การแจกแจงแบบเบ้และค่าเป็นบวกเสมอ}: ตัวแปรเวลา $T \ge 0$ มีลักษณะการแจกแจงแบบเบ้ขวา (Positive Skewness) ซึ่งละเมิดข้อสมมติฐาน Normal Distribution ของ Ordinary Least Squares (OLS)
    \item \textbf{Logistic Regression ขาดมิติของเวลา}: Logistic Regression สนใจเพียงแค่ว่าเกิดเหตุการณ์หรือไม่ ณ จุดสิ้นสุดการทดลอง แต่ไม่สามารถตอบได้ว่าลูกค้ายกเลิกบริการที่เวลาใด หรือความน่าจะเป็นที่จะอยู่รอดผ่านปีที่ 1 หรือปีที่ 2 เป็นเท่าใด
\end{itemize}

\subsection*{2.2 นิยามทางคณิตศาสตร์ของฟังก์ชันการอยู่รอดและความเสี่ยง\\ \normalsize(Mathematical Definitions of Survival \& Hazard Functions)}
กำหนดให้ $T \ge 0$ เป็นตัวแปรสุ่มชนิดต่อเนื่องแทนระยะเวลาจนกว่าจะเกิดเหตุการณ์ โดยมีฟังก์ชันความหนาแน่นความน่าจะเป็น $f(t)$ และฟังก์ชันการแจกแจงสะสม $F(t) = P(T < t)$

\subsubsection*{1. ฟังก์ชันการอยู่รอด (Survival Function: $S(t)$)}
ฟังก์ชัน $S(t)$ คือความน่าจะเป็นที่ลูกค้ารายหนึ่งจะยังคงใช้บริการอยู่เกินกว่าเวลา $t$
\begin{equation}
    S(t) = P(T \ge t) = 1 - F(t) = \int_t^\infty f(u)\,du
\end{equation}
โดยมีคุณสมบัติทางคณิตศาสตร์คือ: $S(0) = 1$ (ณ จุดเริ่มต้น ลูกค้าทุกคนยังอยู่), $S(\infty) = 0$ (เมื่อเวลาผ่านไปไม่สิ้นสุด ทุกคนจะเกิดเหตุการณ์) และ $S(t)$ เป็นฟังก์ชันลดลงทางเดียวหรือไม่เพิ่มขึ้น (Monotonically Non-increasing)

\subsubsection*{2. ฟังก์ชันความเสี่ยง (Hazard Function: $\lambda(t)$ หรือ $h(t)$)}
ฟังก์ชัน $\lambda(t)$ แทนอัตราความเสี่ยงขณะใดขณะหนึ่ง (Instantaneous Failure Rate) ที่ลูกค้าซึ่งรอดชีวิตมาจนถึงเวลา $t$ จะยกเลิกบริการในช่วงเวลาถัดไปสั้น ๆ $\Delta t$
\begin{equation}
    \lambda(t) = \lim_{\Delta t \to 0} \frac{P(t \le T < t + \Delta t \mid T \ge t)}{\Delta t} = \frac{f(t)}{S(t)} = -\frac{d}{dt} \ln S(t)
\end{equation}

\subsubsection*{3. ฟังก์ชันความเสี่ยงสะสม (Cumulative Hazard Function: $\Lambda(t)$)}
คือผลรวมสะสมของความเสี่ยงตั้งแต่จุดเริ่มต้นจนถึงเวลา $t$
\begin{equation}
    \Lambda(t) = \int_0^t \lambda(u)\,du = -\ln S(t) \iff S(t) = \exp\left(-\Lambda(t)\right)
\end{equation}

\subsection*{2.3 ตัวประมาณค่า Kaplan-Meier\\ \normalsize(Kaplan-Meier Non-Parametric Estimator)}
ในการประมาณค่า $S(t)$ แบบนอนพารามิเตอร์โดยไม่สมมติรูปแบบการแจกแจงล่วงหน้า Kaplan \& Meier (1958) ได้พัฒนาสูตรผลคูณสะสม (Product-Limit Formula):
\begin{equation}
    \hat{S}(t) = \prod_{t_i \le t} \left(1 - \frac{d_i}{n_i}\right)
\end{equation}
เมื่อ $t_1 < t_2 < \dots < t_k$ คือลำดับเวลาที่เกิดเหตุการณ์ยกเลิกบริการจริง, $d_i$ คือจำนวนลูกค้าที่เกิดเหตุการณ์ยกเลิกบริการ ณ เวลา $t_i$ และ $n_i$ คือจำนวนลูกค้าที่ยังคงมีความเสี่ยง (Risk Set) ณ จุดเวลาก่อน $t_i$

ช่วงความเชื่อมั่น 95\% คำนวณผ่านความแปรปรวนของ Greenwood (\textbf{Greenwood's Formula}):
\begin{equation}
    \widehat{\text{Var}}(\hat{S}(t)) = \hat{S}(t)^2 \sum_{t_i \le t} \frac{d_i}{n_i(n_i - d_i)}
\end{equation}
\begin{equation}
    95\% \text{ Confidence Interval: } \hat{S}(t) \pm 1.96 \cdot \sqrt{\widehat{\text{Var}}(\hat{S}(t))}
\end{equation}

\subsection*{2.4 การทดสอบสมมติฐาน Log-Rank Test\\ \normalsize(Log-Rank Hypothesis Testing)}
เพื่อทดสอบว่าฟังก์ชันการอยู่รอดของกลุ่มลูกค้า $K$ กลุ่มมีความแตกต่างกันอย่างมีนัยสำคัญหรือไม่ กำหนดสมมติฐาน:
\begin{align*}
    H_0 &: S_1(t) = S_2(t) = \dots = S_K(t) \quad (\text{ทุกกลุ่มมีวิถีการอยู่รอดเหมือนกัน}) \\
    H_1 &: \text{มีอย่างน้อยหนึ่งกลุ่มที่มีฟังก์ชันการอยู่รอดแตกต่างออกไป}
\end{align*}
สถิติทดสอบ Log-Rank ($\chi^2$) เปรียบเทียบผลต่างระหว่างจำนวนเหตุการณ์ที่สังเกตได้จริง ($O_j$) กับค่าคาดหมายภายใต้ $H_0$ ($E_j$):
\begin{equation}
    \chi^2 = \sum_{j=1}^K \frac{(O_j - E_j)^2}{V_j} \sim \chi^2(K - 1)
\end{equation}

\section*{3. โครงสร้างชุดข้อมูลและการเตรียมตัวแปร\\ \large(Dataset Overview \& Preprocessing)}

ชุดข้อมูล \codevar{Telco-Churn.csv} มีโครงสร้างขนาด $7,043$ แถว และ $21$ คอลัมน์ โดยประกอบด้วยตัวแปรที่นำมาใช้ในการวิเคราะห์ดังนี้:
\begin{itemize}[leftmargin=18pt, itemsep=2pt]
    \item \textbf{Duration Variable ($T$)}: คอลัมน์ \codevar{tenure} แทนระยะเวลา (เดือน) ที่ลูกค้าผูกพันกับผู้ให้บริการ ($0 \le T \le 72$)
    \item \textbf{Event Indicator ($E$)}: คอลัมน์ \codevar{Churn} แปลงเป็น $E = 1$ เมื่อ Churn = `Yes' (เกิดเหตุการณ์) และ $E = 0$ เมื่อ Churn = `No' (ถูก Right-Censored)
    \item \textbf{Segmentation Covariate}: คอลัมน์ \codevar{PaymentMethod} แบ่งเป็น 4 กลุ่ม ได้แก่ Electronic check, Mailed check, Bank transfer (automatic) และ Credit card (automatic)
\end{itemize}

\begin{table}[H]
\centering
\caption{สถิติเชิงพรรณนาภาพรวมของชุดข้อมูล Telco Churn ($N = 7,043$)}
\resizebox{0.95\textwidth}{!}{%
\begin{tabular}{l c c c}
\toprule
\textbf{กลุ่มสถานะลูกค้า (Status)} & \textbf{จำนวน (Observations)} & \textbf{สัดส่วน (\%)} & \textbf{ค่าเฉลี่ย Tenure (เดือน)} \\
\midrule
ลูกค้าที่ยกเลิกบริการจริง (Churned Events: $E=1$) & 1,869 & 26.54\% & 17.98 \\
ลูกค้าที่ยังคงใช้บริการอยู่ (Censored: $E=0$) & 5,174 & 73.46\% & 37.57 \\
\midrule
\textbf{รวมประชากรทั้งหมด (Total Cohort)} & \textbf{7,043} & \textbf{100.00\%} & \textbf{32.37} \\
\bottomrule
\end{tabular}%
}
\end{table}

\section*{4. ผลลัพธ์การวิเคราะห์และกราฟสถิติ\\ \large(Experimental Results \& Diagnostic Plots)}

\subsection*{4.1 เส้นโค้งการอยู่รอดของประชากรโดยรวม\\ \normalsize(Overall Survival Curve)}
จากการฟิตโมเดล Kaplan-Meier บนกลุ่มลูกค้าทั้งหมด $7,043$ ราย ได้เส้นโค้งการอยู่รอดดังแสดงใน Figure 1

\begin{figure}[H]
    \centering
    \includegraphics[width=0.82\textwidth]{plot_1_overall_km.pdf}
    \caption{\textbf{Overall Customer Kaplan-Meier Survival Curve $\hat{S}(t)$}: เส้นโค้งการอยู่รอดของลูกค้าภาพรวม ($N=7,043$) พร้อมช่วงความเชื่อมั่น 95\% โดยแสดงค่าความน่าจะเป็น ณ ช่วงเวลาสำคัญ $S(12)=83.9\%$, $S(24)=78.4\%$, $S(48)=68.8\%$ และ $S(72)=58.7\%$}
\end{figure}

\noindent
\textbf{การแปลผล Figure 1}: ค่าความน่าจะเป็นในการอยู่รอดของลูกค้าภาพรวมลดลงอย่างต่อเนื่องแต่คงตัวได้ดี โดยมีลูกค้าถึง $83.9\%$ ที่อยู่รอดเกิน 1 ปี และ $58.7\%$ ที่ยังคงใช้บริการอยู่จนถึงสิ้นสุดการศึกษาที่ 6 ปี (72 เดือน) เนื่องจากมีผู้รอดชีวิตเกิน $50\%$ ตลอด 72 เดือน ค่ามัธยฐานการอยู่รอด (Median Survival Time) จึงมากกว่า 72 เดือน (Not reached within study window)

\subsection*{4.2 เปรียบเทียบเส้นโค้งการอยู่รอดตามช่องทางการชำระเงิน\\ \normalsize(Survival Curves by Payment Method)}
เมื่อจำแนกลูกค้าออกเป็น 4 กลุ่มตามช่องทางการชำระเงิน ได้เส้นโค้งเปรียบเทียบดังแสดงใน Figure 2 และ Figure 3

\begin{figure}[H]
    \centering
    \includegraphics[width=0.85\textwidth]{plot_2_km_by_payment.pdf}
    \caption{\textbf{Comparative Kaplan-Meier Survival Curves by Payment Method}: เปรียบเทียบเส้นโค้งการอยู่รอดระหว่าง 4 ช่องทางชำระเงิน พบว่ากลุ่ม Electronic check มีอัตราการลดลงของความน่าจะเป็นในการอยู่รอดรวดเร็วและรุนแรงที่สุด โดยตัดผ่านเส้นมัธยฐานที่ $t=47.0$ เดือน}
\end{figure}

\begin{figure}[H]
    \centering
    \includegraphics[width=0.92\textwidth]{plot_3_km_grid_by_payment.pdf}
    \caption{\textbf{Individual Kaplan-Meier Facet Plots with 95\% CI}: กราฟแยกย่อย $2 \times 2$ พร้อมแถบช่วงความเชื่อมั่น 95\% แบบ Greenwood แสดงให้เห็นความแม่นยำสูงเนื่องจากขนาดกลุ่มตัวอย่างที่ใหญ่ ($n \ge 1,522$ ในทุกกลุ่ม)}
\end{figure}

\subsection*{4.3 การวิเคราะห์ความเสี่ยงสะสม\\ \normalsize(Cumulative Hazard Function Analysis)}
ฟังก์ชันความเสี่ยงสะสม $\Lambda(t) = -\ln S(t)$ แสดงอัตราการสะสมของความเสี่ยงในการยกเลิกบริการตามระยะเวลาการใช้งาน ดังแสดงใน Figure 4

\begin{figure}[H]
    \centering
    \includegraphics[width=0.85\textwidth]{plot_4_cumulative_hazard.pdf}
    \caption{\textbf{Cumulative Hazard Curves $\Lambda(t)$ by Payment Method}: เส้นความเสี่ยงสะสมของกลุ่ม Electronic check มีความชันสูงมาก ($\Lambda(72) > 1.2$) บ่งชี้ว่าลูกค้ารายใหม่ที่เลือกจ่ายด้วยช่องทางนี้มีความเสี่ยงในการยกเลิกบริการสูงตลอดทุกช่วงเวลา ต่างจากระบบอัตโนมัติที่มีความเสี่ยงสะสมคงที่ในระดับต่ำ ($\Lambda(72) \approx 0.28 - 0.30$)}
\end{figure}

\subsection*{4.4 การแจกแจงอัตรา Churn และระยะเวลา Tenure\\ \normalsize(Churn Rate \& Tenure Distribution)}
Figure 5 แสดงอัตรา Churn (\%) เปรียบเทียบกับกล่องการแจกแจงระยะเวลา Tenure ของลูกค้าแต่ละกลุ่ม

\begin{figure}[H]
    \centering
    \includegraphics[width=0.92\textwidth]{plot_5_churn_distribution.pdf}
    \caption{\textbf{Churn Rate \& Tenure Distribution by Payment Method}: ซ้าย: อัตรา Churn ของ Electronic check สูงถึง $45.3\%$ ขวา: แผนภาพ Boxplot แสดงค่ามัธยฐาน Tenure ของกลุ่มระบบตัดเงินอัตโนมัติที่สูงถึงประมาณ $43$ เดือน}
\end{figure}

\subsection*{4.5 ตารางสรุปผลลัพธ์เชิงปริมาณ\\ \normalsize(Quantitative Summary of Results)}
\begin{table}[H]
\centering
\caption{สรุปผลการวิเคราะห์การอยู่รอดเชิงปริมาณ จำแนกตามช่องทางการชำระเงิน}
\resizebox{\textwidth}{!}{%
\begin{tabular}{l c c c c c c c c c c c}
\toprule
\textbf{Payment Method} & \textbf{Sample $n$} & \textbf{Churn $d$} & \textbf{Censor $c$} & \textbf{Churn Rate} & \textbf{Median Surv.} & \textbf{Mean Tenure} & \textbf{$S(12)$} & \textbf{$S(24)$} & \textbf{$S(36)$} & \textbf{$S(48)$} & \textbf{$S(72)$} \\
\midrule
\textbf{Electronic check} & 2,365 & 1,071 & 1,294 & 45.3\% & \textbf{47.0 mo.} & 25.2 mo. & 0.727 & 0.624 & 0.553 & 0.493 & 0.295 \\
\textbf{Mailed check} & 1,612 & 308 & 1,304 & 19.1\% & > 72 mo. & 21.8 mo. & 0.828 & 0.797 & 0.769 & 0.740 & 0.726 \\
\textbf{Bank transfer (auto)} & 1,544 & 258 & 1,286 & 16.7\% & > 72 mo. & 43.7 mo. & 0.931 & 0.902 & 0.878 & 0.847 & 0.745 \\
\textbf{Credit card (auto)} & 1,522 & 232 & 1,290 & 15.2\% & > 72 mo. & 43.3 mo. & 0.945 & 0.913 & 0.887 & 0.855 & 0.758 \\
\midrule
\textbf{Overall Population} & \textbf{7,043} & \textbf{1,869} & \textbf{5,174} & \textbf{26.5\%} & \textbf{> 72 mo.} & \textbf{32.4 mo.} & \textbf{0.839} & \textbf{0.784} & \textbf{0.738} & \textbf{0.688} & \textbf{0.587} \\
\bottomrule
\end{tabular}%
}
\end{table}

\subsection*{4.6 ผลการทดสอบทางสถิติ Log-Rank Test\\ \normalsize(Log-Rank Statistical Tests)}
\begin{table}[H]
\centering
\caption{ผลการทดสอบความแตกต่างของเส้นโค้งการอยู่รอด (Multivariate \& Pairwise Log-Rank Tests)}
\resizebox{\textwidth}{!}{%
\begin{tabular}{l c c c}
\toprule
\textbf{การเปรียบเทียบ (Comparison)} & \textbf{Chi-Square ($\chi^2$)} & \textbf{$p$-value} & \textbf{ข้อสรุปทางสถิติ ($\alpha=0.001$)} \\
\midrule
\textbf{Multivariate (All 4 Groups)} & \textbf{865.2395} & \textbf{$3.07 \times 10^{-187}$} & \textbf{ปฏิเสธ $H_0$ (แตกต่างกันอย่างมีนัยสำคัญยิ่ง)} \\
\midrule
Bank transfer (auto) vs Credit card (auto) & 0.8683 & $0.3514$ & ไม่แตกต่างกันทางสถิติ ($H_0$ ถูกยอมรับ) \\
Electronic check vs Bank transfer (auto) & 510.0351 & $< 10^{-100}$ & แตกต่างกันอย่างมีนัยสำคัญยิ่ง \\
Electronic check vs Credit card (auto) & 539.7402 & $< 10^{-100}$ & แตกต่างกันอย่างมีนัยสำคัญยิ่ง \\
Electronic check vs Mailed check & 152.4555 & $< 10^{-30}$ & แตกต่างกันอย่างมีนัยสำคัญยิ่ง \\
Mailed check vs Bank transfer (auto) & 51.0719 & $< 10^{-11}$ & แตกต่างกันอย่างมีนัยสำคัญยิ่ง \\
Mailed check vs Credit card (auto) & 64.8152 & $< 10^{-14}$ & แตกต่างกันอย่างมีนัยสำคัญยิ่ง \\
\bottomrule
\end{tabular}%
}
\end{table}

\section*{5. การอภิปรายผลเชิงลึกและข้อมูลเชิงธุรกิจ\\ \large(In-depth Discussion \& Business Insights)}

\subsection*{5.1 การวิเคราะห์กลุ่มความเสี่ยงสูง: Electronic Check\\ \normalsize(High-Risk Segment Analysis)}
ลูกค้ากลุ่มที่ชำระผ่าน \textbf{Electronic check} เป็นกลุ่มที่มีขนาดใหญ่ที่สุด ($n = 2,365$ หรือ $33.58\%$ ของประชากร) แต่กลับมีอัตราการยกเลิกบริการสูงที่สุดถึง \textbf{$45.3\%$} ($1,071$ ราย) และมี \textbf{Median Survival Time เท่ากับ $47.0$ เดือน} 
\begin{itemize}[leftmargin=18pt, itemsep=2pt]
    \item \textbf{พฤติกรรมและความเสียดทานในการชำระ (Payment Friction)}: การจ่ายผ่านเช็คอิเล็กทรอนิกส์มักเป็นการจ่ายแบบ Manual ในแต่ละรอบบิล ทำให้ลูกค้าต้องเผชิญกับ ``จุดตัดสินใจ (Decision Point)'' ทุก ๆ เดือนว่าจะชำระต่อหรือยกเลิกบริการ
    \item \textbf{ความสัมพันธ์กับสัญญาแบบ Month-to-Month}: กลุ่มลูกค้าเช็คอิเล็กทรอนิกส์ส่วนใหญ่มีสัดส่วนสัญญาแบบรายเดือนสูง ทำให้ความยืดหยุ่นในการยกเลิกบริการทำได้ง่ายกว่ากลุ่มที่มีสัญญาระยะยาว
\end{itemize}

\subsection*{5.2 ประสิทธิภาพของระบบชำระเงินอัตโนมัติ\\ \normalsize(Automatic Payment Channels)}
กลุ่มลูกค้าที่ใช้ \textbf{Bank transfer (automatic)} และ \textbf{Credit card (automatic)} มีคุณลักษณะการอยู่รอดที่ดีเยี่ยมและใกล้เคียงกันมาก โดยผลการทดสอบ Pairwise Log-Rank ยืนยันว่าไม่มีความแตกต่างกันอย่างมีนัยสำคัญทางสถิติ ($p = 0.3514$):
\begin{itemize}[leftmargin=18pt, itemsep=2pt]
    \item อัตรา Churn ต่ำมากเพียง $15.2\% - 16.7\%$ 
    \item ค่าความน่าจะเป็นในการอยู่รอดเกิน 2 ปี ($S(24)$) สูงเกิน $90\%$ ($90.2\% - 91.3\%$)
    \item \textbf{กลไกเชิงพฤติกรรม}: ระบบตัดเงินอัตโนมัติช่วยขจัดความตระหนักรู้ถึงค่าใช้จ่ายในแต่ละรอบบิล (Pain of Paying) และลดความยุ่งยากในการทำธุรกรรม ส่งผลให้ลูกค้าคงสถานะการใช้บริการอย่างยาวนาน
\end{itemize}

\subsection*{5.3 พฤติกรรมเฉพาะตัวของกลุ่ม Mailed Check\\ \normalsize(Mailed Check Loyalty Dynamics)}
กลุ่ม \textbf{Mailed check} มีอัตรา Churn รวมค่อนข้างต่ำ ($19.1\%$) และมีค่า $S(72) = 72.6\%$ แต่มีค่าเฉลี่ย Tenure สั้น ($21.8$ เดือน):
\begin{itemize}[leftmargin=18pt, itemsep=2pt]
    \item ในช่วง 12 เดือนแรก มีอัตราการลดลงของเส้นการอยู่รอดที่ค่อนข้างชัน ($S(12) = 82.8\%$) ซึ่งสะท้อนถึงกลุ่มผู้ทดลองใช้บริการระยะสั้น
    \item แต่เมื่อลูกค้ากลุ่มนี้ผ่านช่วง 24 เดือนแรกไปได้ เส้นโค้งการอยู่รอดจะแบนราบลงอย่างเห็นได้ชัด ($S(72) = 72.6\%$) แสดงว่าลูกค้าที่เหลืออยู่คือกลุ่มผู้ใช้ที่มีความภักดีสูง (Brand Loyalists) เช่น กลุ่มผู้สูงอายุหรือลูกค้าดั้งเดิม
\end{itemize}

\section*{6. ข้อเสนอแนะเชิงกลยุทธ์\\ \large(Strategic Recommendations)}

\begin{ansbox}
\textbf{ข้อเสนอแนะเชิงกลยุทธ์สำหรับฝ่ายบริหารและฝ่ายการตลาด:}
\begin{enumerate}[leftmargin=16pt, itemsep=2pt]
    \item \textbf{แคมเปญกระตุ้นการเปลี่ยนผ่านสู่ Auto-Pay (Incentivize Auto-Pay Migration)}: มอบส่วนลดพิเศษ เช่น ส่วนลดค่าบริการ \$5 ต่อเดือน เป็นเวลา 6 เดือน สำหรับลูกค้าที่เปลี่ยนช่องทางจาก Electronic check หรือ Mailed check มาเป็นระบบตัดบัตรเครดิตหรือบัญชีธนาคารอัตโนมัติ ซึ่งจะช่วยลด Churn Rate ได้ทันทีกว่า 28--30\%
    \item \textbf{การแทรกแซงเชิงรุกในช่วง Tenure 0--12 เดือน (Early Onboarding Intervention)}: จากกราฟ Hazard อัตราเสี่ยงต่อการยกเลิกจะสะสมตัวเร็วที่สุดในช่วงปีแรก ควรจัดทีม Customer Success ติดตามความพึงพอใจและเสนอสัญญารายปีพร้อมข้อเสนอพิเศษก่อนที่ลูกค้าจะเข้าสู่เดือนที่ 12
    \item \textbf{ระบบแจ้งเตือนและชำระเงินแบบ One-Click Billing}: ปรับปรุงกระบวนการชำระเงินของ Electronic check ให้มีความสะดวกรวดเร็ว ลดความยุ่งยากของขั้นตอนทางธุรกรรมเพื่อลดแรงจูงใจในการยกเลิกบริการ
\end{enumerate}
\end{ansbox}

\section*{7. สรุปผลการดำเนินงาน\\ \large(Conclusion)}

การวิเคราะห์การอยู่รอดด้วยแบบจำลอง Kaplan-Meier บนชุดข้อมูล Telco Churn ช่วยให้เราเข้าใจพลวัตของระยะเวลาในการคงอยู่ของลูกค้าได้อย่างแม่นยำเหนือกว่าแบบจำลองการถดถอยทั่วไป ผลการทดสอบทางสถิติ Log-Rank Test ($\chi^2 = 865.24, p = 3.07 \times 10^{-187}$) ชี้ชัดว่าช่องทางการชำระเงินเป็นปัจจัยที่มีอิทธิพลต่อความอยู่รอดของลูกค้าอย่างมีนัยสำคัญยิ่ง โดยกลุ่มระบบตัดเงินอัตโนมัติมีประสิทธิภาพในการรักษาฐานลูกค้าสูงสุด ($S(24) > 90\%$) ในขณะที่กลุ่ม Electronic check เป็นกลุ่มความเสี่ยงสูงที่ต้องได้รับการปรับปรุงเชิงกลยุทธ์อย่างเร่งด่วน

\needspace{8\baselineskip}
\section*{สรุปผลลัพธ์ (Summary of Results)}

\begin{table}[H]
\centering
\resizebox{\textwidth}{!}{%
\begin{tabular}{l l}
\toprule
\textbf{หัวข้อ} & \textbf{ผลลัพธ์และข้อสรุป} \\
\midrule
\textbf{ชุดข้อมูล} & Telco Customer Churn ($N = 7,043$): Churn $= 1,869$ ($26.54\%$), Censored $= 5,174$ ($73.46\%$) \\[4pt]
\textbf{Overall Survival} & $S(12) = 0.839, \quad S(24) = 0.784, \quad S(48) = 0.688, \quad S(72) = 0.587$ \\
                          & มัธยฐานการอยู่รอดภาพรวม $> 72$ เดือน (อัตราการรอดชีวิตเกิน $50\%$ ตลอดช่วงศึกษา) \\[4pt]
\textbf{กลุ่มความเสี่ยงสูงสุด} & \textbf{Electronic check} ($n = 2,365$): Churn Rate $= \mathbf{45.3\%}$, Median Survival $= \mathbf{47.0}$ เดือน \\
                             & $S(12) = 0.727 \rightarrow S(24) = 0.624 \rightarrow S(72) = 0.295$ \\[4pt]
\textbf{กลุ่มความเสี่ยงต่ำสุด} & \textbf{Credit card (auto)} ($15.2\%$) \& \textbf{Bank transfer (auto)} ($16.7\%$): $S(24) > 90\%$ \\
                             & มัธยฐานการอยู่รอด $> 72$ เดือน และ Mean Tenure $> 43$ เดือน \\[4pt]
\textbf{Log-Rank Test}       & Multivariate $\chi^2 = \mathbf{865.2395}, \quad p = \mathbf{3.07 \times 10^{-187}}$ (ปฏิเสธ $H_0$) \\
                             & Pairwise: Auto channels ไม่ต่างกัน ($p=0.3514$), Electronic check แย่กว่าทุกกลุ่ม ($p<10^{-30}$) \\[4pt]
\textbf{ข้อเสนอแนะเชิงกลยุทธ์} & รณรงค์เปลี่ยนเป็น Auto-Pay (ลด Churn ทันที $28\%$) และติดตามลูกค้าใหม่ในช่วง $0-12$ เดือน \\
\bottomrule
\end{tabular}%
}
\end{table}

\clearpage
\appendix
\section*{Appendix: Full Jupyter Notebook Code \& Output}

รายละเอียดต่อไปนี้คือโค้ด Python ที่ใช้แก้ปัญหาและวิเคราะห์ชุดข้อมูล Telco Churn พร้อมผลลัพธ์และกราฟทั้งหมด ซึ่งได้รันจริงจาก Jupyter Notebook (\texttt{Assignment\_3\_Regression.ipynb})

% Auto-generated notebook appendix (Assignment 3 - Telco Churn Activity)
\begin{tcolorbox}[colback=blue!5!white,colframe=blue!75!black,title=\textbf{Jupyter Notebook: Assignment\_3\_Regression.ipynb (Telco Churn Survival Analysis)}]
โค้ด ผลลัพธ์ และการพล็อตภาพทั้งหมดด้านล่างเป็นการรันจริงจาก Jupyter Notebook
\end{tcolorbox}

\needspace{5\baselineskip}
\vspace{0.8em}
\needspace{5\baselineskip}
\noindent\textbf{In [1]}:
\begin{lstlisting}[style=pycode]
import os
import numpy as np
import pandas as pd
import matplotlib.pyplot as plt
import matplotlib.font_manager as fm
from lifelines import KaplanMeierFitter, NelsonAalenFitter
from lifelines.statistics import multivariate_logrank_test, pairwise_logrank_test

# --- Robust Sarabun Font Setup ---
font_dir = os.path.join(os.environ.get('LOCALAPPDATA', ''), 'Microsoft', 'Windows', 'Fonts')
sarabun_r_path = os.path.join(font_dir, 'Sarabun-Regular.ttf')
sarabun_b_path = os.path.join(font_dir, 'Sarabun-Bold.ttf')

if os.path.exists(sarabun_r_path):
    fm.fontManager.addfont(sarabun_r_path)
if os.path.exists(sarabun_b_path):
    fm.fontManager.addfont(sarabun_b_path)

plt.rcParams['font.family'] = 'Sarabun'
plt.rcParams['font.size'] = 12
plt.rcParams['axes.unicode_minus'] = False

# --- Data Loading ---
df = pd.read_csv('Telco-Churn.csv')
T = df['tenure']
E = (df['Churn'] == 'Yes').astype(int)

print(f"Telco Churn Dataset Loaded Successfully. Total Observations N = {len(df):,}")
print(f"Observed Churn Events (E=1): {E.sum():,} ({E.mean()*100:.2f}%)")
print(f"Right-Censored Retained Customers (E=0): {(1-E).sum():,} ({(1-E).mean()*100:.2f}%)")

df[['customerID', 'tenure', 'PaymentMethod', 'Contract', 'MonthlyCharges', 'TotalCharges', 'Churn']].head()
\end{lstlisting}
\smallskip
\needspace{5\baselineskip}
\noindent\textbf{Out[1]}:
\begin{lstlisting}[style=output]
Telco Churn Dataset Loaded Successfully. Total Observations N = 7,043
Observed Churn Events (E=1): 1,869 (26.54%)
Right-Censored Retained Customers (E=0): 5,174 (73.46%)
   customerID  tenure  ... TotalCharges Churn
0  7590-VHVEG       1  ...        29.85    No
1  5575-GNVDE      34  ...       1889.5    No
2  3668-QPYBK       2  ...       108.15   Yes
3  7795-CFOCW      45  ...      1840.75    No
4  9237-HQITU       2  ...       151.65   Yes

[5 rows x 7 columns]
\end{lstlisting}
\needspace{10\baselineskip}
\vspace{0.8em}
\begin{tcolorbox}[colback=gray!5!white,colframe=gray!50!black,boxrule=0.5pt,arc=2pt,left=4pt,right=4pt,top=4pt,bottom=4pt,breakable]
\textbf{\large 2. Overall Kaplan-Meier Survival Curve}\par\smallskip
We fit the Kaplan-Meier survival estimator on the full cohort ($N=7,043$) to establish the baseline population survival trajectory $\hat{S}(t)$ over tenure.\par
\end{tcolorbox}
\needspace{5\baselineskip}
\vspace{0.8em}
\needspace{5\baselineskip}
\noindent\textbf{In [2]}:
\begin{lstlisting}[style=pycode]
# --- Overall Survival Estimation ---
kmf_overall = KaplanMeierFitter()
kmf_overall.fit(T, event_observed=E, label='All Customers (N = 7,043)')

fig, ax = plt.subplots(figsize=(8.5, 5))
kmf_overall.plot_survival_function(ax=ax, color='#1e3a8a', linewidth=2.5, ci_show=True, ci_alpha=0.2)

# Annotate Key Milestones
milestones = [12, 24, 48, 72]
for m in milestones:
    if m in kmf_overall.survival_function_.index:
        surv_val = kmf_overall.predict(m)
        ax.plot(m, surv_val, marker='o', markersize=6, color='#dc2626')
        ax.annotate(f"{surv_val:.1%}\n(t={m})", (m, surv_val), textcoords="offset points",
                    xytext=(0, 10), ha='center', fontsize=9, fontweight='bold')

ax.set_title('Overall Customer Kaplan-Meier Survival Curve $\\hat{S}(t)$', fontsize=14, fontweight='bold', pad=12)
ax.set_xlabel('Tenure (Months)', fontsize=12)
ax.set_ylabel('Survival Probability $S(t)$', fontsize=12)
ax.set_ylim(0, 1.05)
ax.set_xlim(0, 75)
ax.grid(True, linestyle='--', alpha=0.5)
plt.tight_layout()
plt.savefig('plot_1_overall_km.pdf', bbox_inches='tight')
plt.show()
\end{lstlisting}
\noindent\begin{minipage}{\linewidth}
\centering
\vspace{0.4em}
\includegraphics[width=0.96\linewidth]{plot_1_overall_km.pdf}
\par\vspace{0.6em}
\begin{tcolorbox}[colback=gray!5!white,colframe=gray!50!black,boxrule=0.5pt,arc=2pt,left=6pt,right=6pt,top=5pt,bottom=5pt,unbreakable]
\textbf{Figure 1 (Overall Survival Curve)}: Kaplan-Meier non-parametric survival function $\hat{S}(t)$ for the full customer cohort ($N = 7,043$) with 95\% Greenwood confidence interval. Milestone annotations illustrate that 83.9\% of customers survive past 1 year ($t=12$), 78.4\% survive past 2 years ($t=24$), and 58.7\% remain active at 6 years ($t=72$). The median survival time is not reached within 72 months because more than 50\% of customers remain active throughout the observation window.\par
\end{tcolorbox}
\vspace{1.0em}
\end{minipage}
\needspace{10\baselineskip}
\vspace{0.8em}
\begin{tcolorbox}[colback=gray!5!white,colframe=gray!50!black,boxrule=0.5pt,arc=2pt,left=4pt,right=4pt,top=4pt,bottom=4pt,breakable]
\textbf{\large 3. Comparative Survival Analysis by Payment Method}\par\smallskip
We segment customers into 4 distinct groups based on \texttt{PaymentMethod} to compare their survival trajectories and hazard characteristics.\par
\end{tcolorbox}
\needspace{5\baselineskip}
\vspace{0.8em}
\needspace{5\baselineskip}
\noindent\textbf{In [3]}:
\begin{lstlisting}[style=pycode]
methods = ['Electronic check', 'Mailed check', 'Bank transfer (automatic)', 'Credit card (automatic)']
palette = {
    'Electronic check': '#dc2626',
    'Mailed check': '#ea580c',
    'Bank transfer (automatic)': '#2563eb',
    'Credit card (automatic)': '#16a34a'
}

fig, ax = plt.subplots(figsize=(9.5, 5.5))
kmf_dict = {}

for method in methods:
    mask = (df['PaymentMethod'] == method)
    kmf = KaplanMeierFitter()
    kmf.fit(T[mask], event_observed=E[mask], label=method)
    kmf_dict[method] = kmf
    kmf.plot_survival_function(ax=ax, color=palette[method], linewidth=2.2, ci_show=True, ci_alpha=0.12)

ax.set_title('Comparative Kaplan-Meier Survival Curves by Payment Method', fontsize=14, fontweight='bold', pad=12)
ax.set_xlabel('Tenure (Months)', fontsize=12)
ax.set_ylabel('Survival Probability $S(t)$', fontsize=12)
ax.set_ylim(0, 1.05)
ax.set_xlim(0, 75)
ax.axhline(0.5, color='gray', linestyle=':', alpha=0.7, label='Median Survival Threshold ($S(t)=0.5$)')
ax.grid(True, linestyle='--', alpha=0.5)
ax.legend(title='Payment Method', loc='lower left', frameon=True, framealpha=0.9)
plt.tight_layout()
plt.savefig('plot_2_km_by_payment.pdf', bbox_inches='tight')
plt.show()
\end{lstlisting}
\noindent\begin{minipage}{\linewidth}
\centering
\vspace{0.4em}
\includegraphics[width=0.96\linewidth]{plot_2_km_by_payment.pdf}
\par\vspace{0.6em}
\begin{tcolorbox}[colback=gray!5!white,colframe=gray!50!black,boxrule=0.5pt,arc=2pt,left=6pt,right=6pt,top=5pt,bottom=5pt,unbreakable]
\textbf{Figure 2 (Comparative Survival Curves)}: Kaplan-Meier survival curves comparing the four payment methods. Customers utilizing \textbf{Electronic check} exhibit a steep, continuous drop in survival probability, crossing the 50\% median threshold at $t = 47.0$ months. Conversely, automated payment channels (\textbf{Bank transfer} and \textbf{Credit card}) display nearly identical, superior retention curves maintaining $> 90\%$ survival past 2 years and $> 74\%$ at 6 years. \textbf{Mailed check} shows moderate early drop-off followed by long-term stabilization.\par
\end{tcolorbox}
\vspace{1.0em}
\end{minipage}
\needspace{10\baselineskip}
\vspace{0.8em}
\begin{tcolorbox}[colback=gray!5!white,colframe=gray!50!black,boxrule=0.5pt,arc=2pt,left=4pt,right=4pt,top=4pt,bottom=4pt,breakable]
\textbf{\large 4. Individual Facet Plots with Confidence Bands}\par\smallskip
Plotting each payment method separately in a $2 \times 2$ grid to clearly inspect the 95\% Greenwood confidence intervals and group characteristics.\par
\end{tcolorbox}
\needspace{5\baselineskip}
\vspace{0.8em}
\needspace{5\baselineskip}
\noindent\textbf{In [4]}:
\begin{lstlisting}[style=pycode]
fig, axes = plt.subplots(2, 2, figsize=(11, 8), sharex=True, sharey=True)
axes = axes.flatten()

for i, method in enumerate(methods):
    ax = axes[i]
    mask = (df['PaymentMethod'] == method)
    kmf = kmf_dict[method]
    kmf.plot_survival_function(ax=ax, color=palette[method], linewidth=2.2, ci_show=True, ci_alpha=0.25)
    
    n_total = len(T[mask])
    n_churn = E[mask].sum()
    churn_rate = (n_churn / n_total) * 100
    med_val = kmf.median_survival_time_
    med_str = f"{med_val:.1f} mo." if not np.isinf(med_val) else "> 72 mo. (N/A)"
    
    ax.set_title(f"{method}\n(n={n_total:,}, Churn={churn_rate:.1f}%, Median={med_str})", fontsize=11, fontweight='bold')
    ax.set_ylim(0, 1.05)
    ax.set_xlim(0, 75)
    ax.grid(True, linestyle='--', alpha=0.5)
    ax.set_xlabel('Tenure (Months)', fontsize=10)
    ax.set_ylabel('Survival Probability $S(t)$', fontsize=10)
    ax.axhline(0.5, color='gray', linestyle=':', alpha=0.5)

fig.suptitle('Individual Kaplan-Meier Survival Estimates with 95% Confidence Intervals', fontsize=14, fontweight='bold', y=0.98)
plt.tight_layout()
plt.savefig('plot_3_km_grid_by_payment.pdf', bbox_inches='tight')
plt.close()
print("Plot 3 saved to plot_3_km_grid_by_payment.pdf")
\end{lstlisting}
\needspace{5\baselineskip}
\smallskip
\needspace{5\baselineskip}
\noindent\textbf{Out[4]}:
\begin{lstlisting}[style=output]
Plot 3 saved to plot_3_km_grid_by_payment.pdf
\end{lstlisting}
\vspace{0.4em}
\noindent\begin{minipage}{\linewidth}
\centering
\vspace{0.4em}
\includegraphics[width=0.96\linewidth]{plot_3_km_grid_by_payment.pdf}
\par\vspace{0.6em}
\begin{tcolorbox}[colback=gray!5!white,colframe=gray!50!black,boxrule=0.5pt,arc=2pt,left=6pt,right=6pt,top=5pt,bottom=5pt,unbreakable]
\textbf{Figure 3 (Individual Survival Facets with 95\% CI)}: $2 \times 2$ facet grid of Kaplan-Meier survival curves with 95\% confidence bands for each payment method. The tight confidence intervals across all four subplots confirm high statistical precision resulting from large subgroup sample sizes ($n \ge 1,522$ per group).\par
\end{tcolorbox}
\vspace{1.0em}
\end{minipage}
\needspace{10\baselineskip}
\vspace{0.8em}
\begin{tcolorbox}[colback=gray!5!white,colframe=gray!50!black,boxrule=0.5pt,arc=2pt,left=4pt,right=4pt,top=4pt,bottom=4pt,breakable]
\textbf{\large 5. Cumulative Hazard Function Analysis}\par\smallskip
The cumulative hazard function $\Lambda(t) = \int_0^t \lambda(u)\,du = -\ln S(t)$ models the accumulated risk of customer churn over time.\par
\end{tcolorbox}
\needspace{5\baselineskip}
\vspace{0.8em}
\needspace{5\baselineskip}
\noindent\textbf{In [5]}:
\begin{lstlisting}[style=pycode]
fig, ax = plt.subplots(figsize=(9, 5.5))
naf_dict = {}

for method in methods:
    mask = (df['PaymentMethod'] == method)
    naf = NelsonAalenFitter()
    naf.fit(T[mask], event_observed=E[mask], label=method)
    naf_dict[method] = naf
    naf.plot_cumulative_hazard(ax=ax, color=palette[method], linewidth=2.2, ci_show=True, ci_alpha=0.12)

ax.set_title('Cumulative Hazard $\\Lambda(t)$ by Payment Method', fontsize=14, fontweight='bold', pad=12)
ax.set_xlabel('Tenure (Months)', fontsize=12)
ax.set_ylabel('Cumulative Hazard $\\Lambda(t) = -\\ln S(t)$', fontsize=12)
ax.set_xlim(0, 75)
ax.grid(True, linestyle='--', alpha=0.5)
ax.legend(title='Payment Method', loc='upper left', frameon=True, framealpha=0.9)
plt.tight_layout()
plt.savefig('plot_4_cumulative_hazard.pdf', bbox_inches='tight')
plt.show()
\end{lstlisting}
\noindent\begin{minipage}{\linewidth}
\centering
\vspace{0.4em}
\includegraphics[width=0.96\linewidth]{plot_4_cumulative_hazard.pdf}
\par\vspace{0.6em}
\begin{tcolorbox}[colback=gray!5!white,colframe=gray!50!black,boxrule=0.5pt,arc=2pt,left=6pt,right=6pt,top=5pt,bottom=5pt,unbreakable]
\textbf{Figure 4 (Cumulative Hazard Curves)}: Cumulative hazard function $\Lambda(t) = -\ln S(t)$ across payment methods. Electronic check exhibits a steeply ascending, quasi-linear hazard rate surpassing $\Lambda(72) > 1.2$, indicating severe ongoing churn risk. In contrast, automated methods exhibit flat, gradual hazard accumulation ($\Lambda(72) \approx 0.28$ to $0.30$).\par
\end{tcolorbox}
\vspace{1.0em}
\end{minipage}
\needspace{14\baselineskip}
\vspace{0.8em}
\begin{tcolorbox}[colback=gray!5!white,colframe=gray!50!black,boxrule=0.5pt,arc=2pt,left=4pt,right=4pt,top=4pt,bottom=4pt,breakable]
\textbf{\large 6. Distribution Analysis of Churn Rate \& Tenure}\par\smallskip
Examining the overall churn percentage and the distribution of customer tenure across the four payment methods.\par
\end{tcolorbox}
\needspace{5\baselineskip}
\vspace{0.8em}
\needspace{5\baselineskip}
\noindent\textbf{In [6]}:
\begin{lstlisting}[style=pycode]
fig, (ax1, ax2) = plt.subplots(1, 2, figsize=(12, 5))

# Churn Rate Bar Chart
churn_summary = df.groupby('PaymentMethod')['Churn'].apply(lambda x: (x == 'Yes').mean() * 100).reindex(methods)
bars = ax1.bar(range(len(methods)), churn_summary, color=[palette[m] for m in methods], edgecolor='black', alpha=0.85, width=0.6)
ax1.set_xticks(range(len(methods)))
ax1.set_xticklabels([m.replace(' (automatic)', '\n(auto)') for m in methods], fontsize=10)
ax1.set_ylabel('Churn Rate (%)', fontsize=11)
ax1.set_title('Customer Churn Rate by Payment Method', fontsize=12, fontweight='bold')
ax1.set_ylim(0, 55)
ax1.grid(True, linestyle='--', axis='y', alpha=0.5)
for bar in bars:
    yval = bar.get_height()
    ax1.text(bar.get_x() + bar.get_width()/2.0, yval + 1, f"{yval:.1f}%", ha='center', va='bottom', fontweight='bold', fontsize=10)

# Tenure Boxplot
data_to_plot = [df[df['PaymentMethod'] == m]['tenure'] for m in methods]
bp = ax2.boxplot(data_to_plot, patch_artist=True, tick_labels=[m.replace(' (automatic)', '\n(auto)') for m in methods],
                 medianprops=dict(color='black', linewidth=1.5))
for patch, m in zip(bp['boxes'], methods):
    patch.set_facecolor(palette[m])
    patch.set_alpha(0.6)
ax2.set_ylabel('Tenure (Months)', fontsize=11)
ax2.set_title('Tenure Distribution by Payment Method', fontsize=12, fontweight='bold')
ax2.grid(True, linestyle='--', axis='y', alpha=0.5)

plt.tight_layout()
plt.savefig('plot_5_churn_distribution.pdf', bbox_inches='tight')
plt.show()
\end{lstlisting}
\noindent\begin{minipage}{\linewidth}
\centering
\vspace{0.4em}
\includegraphics[width=0.96\linewidth]{plot_5_churn_distribution.pdf}
\par\vspace{0.6em}
\begin{tcolorbox}[colback=gray!5!white,colframe=gray!50!black,boxrule=0.5pt,arc=2pt,left=6pt,right=6pt,top=5pt,bottom=5pt,unbreakable]
\textbf{Figure 5 (Churn Rate \& Tenure Distribution)}: Left: Churn rates across payment methods highlighting Electronic check ($45.3\%$) vs. Mailed check ($19.1\%$), Bank transfer ($16.7\%$), and Credit card ($15.2\%$). Right: Boxplots of tenure demonstrating substantially higher customer loyalty and tenure in automated payment channels (median tenure $\approx 43$ months) compared to manual payment methods.\par
\end{tcolorbox}
\vspace{1.0em}
\end{minipage}
\needspace{10\baselineskip}
\vspace{0.8em}
\begin{tcolorbox}[colback=gray!5!white,colframe=gray!50!black,boxrule=0.5pt,arc=2pt,left=4pt,right=4pt,top=4pt,bottom=4pt,breakable]
\textbf{\large 7. Statistical Hypothesis Testing \& Quantitative Summary}\par\smallskip
We perform the \textbf{Multivariate Log-Rank Test} to test $H_0: S_1(t) = S_2(t) = S_3(t) = S_4(t)$, followed by pairwise comparisons and milestone survival metrics compilation.\par
\end{tcolorbox}
\needspace{5\baselineskip}
\vspace{0.8em}
\needspace{5\baselineskip}
\noindent\textbf{In [7]}:
\begin{lstlisting}[style=pycode]
# --- Log-Rank Hypothesis Testing ---
logrank_res = multivariate_logrank_test(T, df['PaymentMethod'], E)
print("="*60)
print("MULTIVARIATE LOG-RANK TEST RESULTS")
print("="*60)
print(f"Chi-Square: {logrank_res.test_statistic:.4f}")
print(f"p-value:    {logrank_res.p_value:.4e}")
print(f"df:         {logrank_res.degrees_of_freedom}")
print(f"Conclusion: Reject H0 at alpha = 0.001")

pairwise_res = pairwise_logrank_test(T, df['PaymentMethod'], E)
print("\n" + "="*60)
print("PAIRWISE LOG-RANK TEST SUMMARY")
print("="*60)
for pair, row in pairwise_res.summary.iterrows():
    p_val = row['p']
    p_str = f"{p_val:.4e}" if p_val < 0.001 else f"{p_val:.4f}"
    print(f"{pair[0][:15]:<16} vs {pair[1][:15]:<16} | chi2={row['test_statistic']:>7.2f} | p={p_str}")

# --- Quantitative Summary Table ---
summary_list = []
for m in methods:
    mask = (df['PaymentMethod'] == m)
    kmf = kmf_dict[m]
    n_tot = len(T[mask])
    n_ev = E[mask].sum()
    n_cens = n_tot - n_ev
    ch_rate = (n_ev / n_tot) * 100
    med_s = kmf.median_survival_time_
    med_str = f"{med_s:.1f} mo." if not np.isinf(med_s) else "> 72 mo."
    mean_t = T[mask].mean()
    
    summary_list.append({
        'Payment Method': m,
        'Sample (n)': n_tot,
        'Churned (d)': n_ev,
        'Censored (c)': n_cens,
        'Churn Rate': f"{ch_rate:.1f}%",
        'Median Survival': med_str,
        'Mean Tenure': f"{mean_t:.1f} mo.",
        'S(12 mo.)': f"{kmf.predict(12):.3f}",
        'S(24 mo.)': f"{kmf.predict(24):.3f}",
        'S(36 mo.)': f"{kmf.predict(36):.3f}",
        'S(48 mo.)': f"{kmf.predict(48):.3f}",
        'S(60 mo.)': f"{kmf.predict(60):.3f}",
        'S(72 mo.)': f"{kmf.predict(72):.3f}"
    })

sum_df = pd.DataFrame(summary_list)
print("\n" + "="*60)
print("COMPREHENSIVE SURVIVAL SUMMARY TABLE")
print("="*60)
sum_df
\end{lstlisting}
\smallskip
\newpage
\noindent\textbf{Out[7]}:
\begin{lstlisting}[style=output]
============================================================
MULTIVARIATE LOG-RANK TEST RESULTS
============================================================
Chi-Square: 865.2395
p-value:    3.0665e-187
df:         3
Conclusion: Reject H0 at alpha = 0.001

============================================================
PAIRWISE LOG-RANK TEST SUMMARY
============================================================
Bank transfer (  vs Credit card (au  | chi2=   0.87 | p=0.3514
Bank transfer (  vs Electronic chec  | chi2= 510.04 | p=6.2314e-113
Bank transfer (  vs Mailed check     | chi2=  51.07 | p=8.9046e-13
Credit card (au  vs Electronic chec  | chi2= 539.74 | p=2.1476e-119
Credit card (au  vs Mailed check     | chi2=  64.82 | p=8.2263e-16
Electronic chec  vs Mailed check     | chi2= 152.46 | p=5.0383e-35

============================================================
COMPREHENSIVE SURVIVAL SUMMARY TABLE
============================================================
              Payment Method  Sample (n)  ...  S(60 mo.)  S(72 mo.)
0           Electronic check        2365  ...      0.413      0.295
1               Mailed check        1612  ...      0.732      0.726
2  Bank transfer (automatic)        1544  ...      0.804      0.745
3    Credit card (automatic)        1522  ...      0.831      0.758

[4 rows x 13 columns]
\end{lstlisting}
\needspace{10\baselineskip}
\vspace{0.8em}
\begin{tcolorbox}[colback=gray!5!white,colframe=gray!50!black,boxrule=0.5pt,arc=2pt,left=4pt,right=4pt,top=4pt,bottom=4pt,breakable]
\textbf{\large 8. Summary of Results \& Domain Discussion}\par\smallskip
\smallskip
\textbf{Summary of Quantitative Findings:}\par\smallskip
1. \textbf{Electronic Check Customers (High Risk Channel)}:\par
\textbullet\ Has the largest subgroup ($n = 2,365, 33.6\%$ of total cohort) but suffers an alarmingly high churn rate of \textbf{$45.3\%$} ($1,071$ churn events).\par
\textbullet\ Reaches median survival time at \textbf{$47.0$ months} ($S(47.0) = 0.500$).\par
\textbullet\ Only \textbf{$72.7\%$} survive past 1 year ($t=12$), dropping to \textbf{$62.4\%$} at 2 years ($t=24$) and \textbf{$29.5\%$} at 6 years ($t=72$).\par
\textbullet\ \textbf{Domain Insight:} Electronic check billing requires manual monthly intervention without the automated stickiness of auto-pay, making customers actively reassess and cancel their subscriptions.\par
\smallskip
2. \textbf{Automatic Payment Channels (Bank Transfer \& Credit Card - Low Risk Loyalty Channels)}:\par
\textbullet\ Both automated channels display exceptional, virtually identical survival characteristics ($p = 0.351$, not statistically different from each other).\par
\textbullet\ \textbf{Credit Card (automatic):} Churn rate = \textbf{$15.2\%$}, $S(12) = 0.945$, $S(24) = 0.913$, $S(72) = 0.758$, Mean tenure = $43.3$ months.\par
\textbullet\ \textbf{Bank Transfer (automatic):} Churn rate = \textbf{$16.7\%$}, $S(12) = 0.931$, $S(24) = 0.902$, $S(72) = 0.745$, Mean tenure = $43.7$ months.\par
\textbullet\ Neither method reaches the 50\% survival median within the 72-month study window (Median $> 72$ months).\par
\textbullet\ \textbf{Domain Insight:} Auto-pay minimizes friction during recurring billing cycles, shielding customers from churn decision fatigue.\par
\smallskip
3. \textbf{Mailed Check Customers (Moderate Churn with Long-Term Stabilization)}:\par
\textbullet\ Sample size $n = 1,612$, Churn rate = \textbf{$19.1\%$}, Mean tenure = $21.8$ months.\par
\textbullet\ Exhibits steep initial churn ($S(12) = 0.828, S(24) = 0.797$) during trial onboarding, but the curve flattens significantly after 2 years ($S(72) = 0.726$).\par
\smallskip
\textbf{Strategic Business Recommendations:}\par\smallskip
\textbullet\ \textbf{Incentivize Auto-Pay Adoption:} Implement promotional discounts (\$5/mo discount for 6 months) to convert Electronic check and Mailed check users to Credit card or Bank transfer auto-pay.\par
\textbullet\ \textbf{Early Intervention (Tenure $0 - 12$ months):} Automated onboarding check-ins and proactive customer support during the first 12 months to curb initial hazard spikes.\par
\end{tcolorbox}

\end{document}
